\section{物理层 III：振荡、相位与同步}

文献地图中的同步教材之所以重要，是因为许多转变都耦合于外部时间结构。
一个日程、一个搭档，或者一个制度性时钟，都可能把行为牵引进同步。

\begin{discussion}
同步首先不是关于社会压力，而是关于相位锁定。
如果两个或更多过程反复对齐，目标系统偏离到任意时序的自由度就会减少。
这也正是为什么共享起始时间、重复课程表与共学时段，往往比完全私人决定的开始更容易维持：
耦合减少了时间不确定性。
\end{discussion}

\begin{example}[共学作为相位控制]
两名学生约定每天晚上 8:00 一起开始。
这个安排的价值，不只是因为每个人都害怕让对方失望。
更关键的是，状态转变的时机被外部耦合了。
因此，这一干预作用于相位对齐，而不仅仅作用于动机。
\end{example}